% =====================================================================
%  CHƯƠNG 5. LÝ THUYẾT HỌC TĂNG CƯỜNG THEO GEEKSFORGEEKS
% =====================================================================
\chapter[Lý thuyết học tăng cường theo GeeksforGeeks]{Lý thuyết học tăng cường\\theo GeeksforGeeks}

\section{Bản chất của học tăng cường (Reinforcement Learning -- RL)}
Theo GeeksforGeeks \cite{gfg_rl}, học tăng cường là một nhánh của học máy, trong đó một tác tử tự trị học cách đưa ra
chuỗi quyết định tối ưu bằng cách tương tác với môi trường để tối đa hóa tổng phần thưởng tích lũy. Tác tử không được chỉ
dẫn hành động nào là đúng; nó học qua \textbf{thử và sai} (trial and error), nhận phần thưởng khi làm tốt và hình phạt khi
làm sai. Các thành phần cốt lõi gồm: \textbf{tác tử} (agent), \textbf{môi trường} (environment), \textbf{trạng thái},
\textbf{hành động}, \textbf{phần thưởng}, \textbf{chính sách} (policy) và \textbf{hàm giá trị} (value function).

Điểm khác biệt với hai mô thức học máy còn lại (Hình~\ref{fig:ch5-types}):
\begin{itemize}
  \item \textbf{So với học có giám sát (Supervised):} không có tập dữ liệu gán nhãn đúng/sai từ trước; tác tử phải tự
        khám phá hành vi tối ưu qua tín hiệu thưởng/phạt, và tín hiệu này có thể đến trễ nhiều bước.
  \item \textbf{So với học không giám sát (Unsupervised):} không chỉ tìm cấu trúc ẩn trong dữ liệu mà hướng tới cực đại
        hóa một hàm mục tiêu số học là tổng phần thưởng.
\end{itemize}

\hinhtikz{ch5_learning_types}{Ba mô thức học máy: học có giám sát, học không giám sát và học tăng cường.}{ch5-types}

\section{Bài toán đánh đổi thăm dò (Exploration) và khai thác (Exploitation)}
Tác tử luôn đứng trước một thế khó: \emph{khai thác} hành động đã biết để nhận phần thưởng chắc chắn, hay \emph{thăm dò}
hành động mới để có thể tìm ra phần thưởng lớn hơn. Chỉ khai thác thì dễ mắc kẹt ở lời giải chưa tối ưu; chỉ thăm dò thì
không bao giờ tận dụng được điều đã học. Chiến lược phổ biến nhất là \textbf{$\varepsilon$-greedy}:
\[
  a =
  \begin{cases}
    \text{hành động ngẫu nhiên trong } \mathcal{A}, & \text{với xác suất } \varepsilon \quad \text{(thăm dò)},\\[2pt]
    \arg\max_{a'} Q(s, a'), & \text{với xác suất } 1-\varepsilon \quad \text{(khai thác)}.
  \end{cases}
\]
Trong thực tế $\varepsilon$ được giảm dần theo thời gian: đầu quá trình tác tử chưa biết gì nên cần thăm dò nhiều; khi
Q-Table đã đáng tin, tác tử chuyển sang khai thác. Bài thực hành dùng lịch giảm theo cấp số nhân
$\varepsilon_{k+1} = \max(\varepsilon_{\min},\, \varepsilon_k \cdot \lambda)$ với $\varepsilon_0 = 1$,
$\varepsilon_{\min} = 0{,}01$ và $\lambda = 0{,}998$.

\hinhtikz{ch5_eps_greedy}{Chính sách $\varepsilon$-greedy: phân nhánh thăm dò -- khai thác và lịch giảm dần của
$\varepsilon$.}{ch5-eps}

\section{Hàm giá trị và phương trình tối ưu Bellman}
Hàm giá trị hành động $Q^\pi(s, a)$ biểu thị kỳ vọng tổng phần thưởng chiết khấu khi tác tử thực hiện hành động $a$ tại
trạng thái $s$ rồi sau đó tuân theo chính sách $\pi$:
\[
  Q^\pi(s,a) = \mathbb{E}_\pi\!\left[ G_t \mid S_t = s, A_t = a \right].
\]
Hàm giá trị tối ưu $Q^*$ thỏa \textbf{phương trình tối ưu Bellman} \cite{sutton_barto}:
\[
  Q^*(s,a) = \mathbb{E}\!\left[ R_{t+1} + \gamma \max_{a'} Q^*(S_{t+1}, a') \;\middle|\; S_t = s, A_t = a \right].
\]
Thuật toán \textbf{Q-Learning} xấp xỉ $Q^*$ bằng cách lấy mẫu từ kinh nghiệm thực tế và cập nhật dần theo quy tắc
\[
  Q(s,a) \;\leftarrow\; Q(s,a) + \alpha \Big[ \underbrace{R + \gamma \max_{a'} Q(s', a')}_{\text{mục tiêu TD}} - Q(s,a)
  \Big],
\]
trong đó $\alpha \in (0,1]$ là tốc độ học và $\gamma \in [0,1]$ là hệ số chiết khấu. Phần trong ngoặc vuông là \emph{sai
số TD} (temporal difference). Khi tác tử về trạng thái kết thúc, mục tiêu TD chỉ còn là $R$ vì không còn phần thưởng
tương lai. Q-Learning là thuật toán \emph{off-policy}: nó học chính sách tham lam tối ưu trong khi hành động theo
$\varepsilon$-greedy.

\hinhtikz{ch5_bellman}{Sơ đồ sao lưu (backup diagram) của phép cập nhật Q-Learning.}{ch5-bellman}

\section{Phân loại các trường phái học tăng cường}
Bảng~\ref{tab:ch5-families} và Hình~\ref{fig:ch5-families} tóm tắt bốn trường phái chính.

\begin{table}[!htbp]\centering
\caption{Phân loại các trường phái học tăng cường chính.}\label{tab:ch5-families}
\small
\begin{tabularx}{\linewidth}{L{2.8cm}XL{4.4cm}}
\toprule
\textbf{Trường phái} & \textbf{Đặc trưng kỹ thuật cốt lõi} & \textbf{Thuật toán tiêu biểu} \\
\midrule
Value-Based & Ước lượng hàm giá trị tối ưu $Q^*(s,a)$, chọn hành động tham lam theo giá trị & Q-Learning, SARSA,
Deep Q-Network (DQN) \\
Policy-Based & Tối ưu trực tiếp chính sách $\pi_\theta(a \mid s)$ bằng gradient ascent, không cần Q-Table & REINFORCE,
Policy Gradient, TRPO \\
Actor-Critic & Kết hợp cả hai: Actor sinh hành động, Critic đánh giá giá trị để giảm phương sai & A2C, A3C, PPO \\
Model-Based & Tự xây dựng mô hình giả lập môi trường để lập kế hoạch trước khi hành động & Dyna-Q, World Models,
MuZero \\
\bottomrule
\end{tabularx}
\end{table}

\hinhtikz{ch5_rl_families}{Bốn trường phái học tăng cường và các thuật toán tiêu biểu.}{ch5-families}

GeeksforGeeks \cite{gfg_rl} cũng nêu các ứng dụng điển hình của học tăng cường: robot và điều khiển tự động, trò chơi
(cờ vây, game điện tử), xe tự lái, hệ gợi ý, tài chính định lượng, và tinh chỉnh mô hình ngôn ngữ theo phản hồi con
người (Chương~\ref{ch:giaothoa}).
